\section{هفته‌ی 2: کالبدشکافی یک مدل: تشخیص علامت و منحنی‌های فراموشی}

\begin{quote}
\textbf{فصل 1: مبانی مدل‌سازی شناختی} · درس 1 ساعت · \lr{lab} 2 ساعت
\textbf{مدل‌های اصلی این هفته:} نظریه‌ی تشخیص علامت گاوسی با واریانس برابر
\lr{(SDT)}؛ منحنی‌های فراموشی نمایی و توانی. هر دو در هفته‌های 3 تا 13
دوباره می‌آیند.
\end{quote}

\begin{center}\rule{0.5\linewidth}{0.5pt}\end{center}

\subsection{واژه‌های کلیدی}

واژه‌های مهم این درس با معادل فارسی‌شان (حتی اگر در درس‌های قبل آمده
باشند).

{\def\LTcaptype{none} % do not increment counter
\begin{longtable}[]{@{}
  >{\raggedright\arraybackslash}p{(\linewidth - 4\tabcolsep) * \real{0.3333}}
  >{\raggedright\arraybackslash}p{(\linewidth - 4\tabcolsep) * \real{0.3333}}
  >{\raggedright\arraybackslash}p{(\linewidth - 4\tabcolsep) * \real{0.3333}}@{}}
\toprule\noalign{}
\begin{minipage}[b]{\linewidth}\raggedright
\lr{English}
\end{minipage} & \begin{minipage}[b]{\linewidth}\raggedright
فارسی
\end{minipage} & \begin{minipage}[b]{\linewidth}\raggedright
توضیح کوتاه
\end{minipage} \\
\midrule\noalign{}
\endhead
\bottomrule\noalign{}
\endlastfoot
\lr{forgetting curve} & منحنی فراموشی & کاهش یادآوری با گذشت زمان \\
\lr{signal detection theory (SDT)} & نظریه‌ی تشخیص علامت & مدلی که حساسیت
را از سوگیری پاسخ جدا می‌کند \\
\lr{exponential} / \lr{power function} & تابع نمایی / توانی & دو شکل
رایج برای منحنی فراموشی و یادگیری \\
\lr{parameter} & پارامتر & کمیت قابل تنظیم مدل که بین افراد یا شرایط فرق
می‌کند \\
\lr{deterministic core} & هسته‌ی قطعی & بخش بدون تصادف مدل که پیش‌بینی
میانگین را می‌دهد \\
\lr{noise} & نویز & تغییرات تصادفی در شواهد یا در داده \\
\lr{linking assumption} & فرض پیوند & اینکه کمیت‌های مدل چطور به رفتار
قابل مشاهده وصل می‌شوند \\
\lr{sensitivity (d′)} & حساسیت & فاصله‌ی توزیع علامت و نویز \\
\lr{hit} & اصابت \lr{(hit)} & پاسخ «بله» وقتی علامت وجود دارد \\
\lr{response bias} & سوگیری پاسخ & تمایل به یک پاسخ، مستقل از محرک \\
\lr{binomial distribution} & توزیع دوجمله‌ای & توزیع تعداد موفقیت در
\lr{n} کوشش \\
\lr{noise model} & مدل نویز & توزیع داده حول پیش‌بینی مدل \\
\lr{simulation} & شبیه‌سازی & تولید داده‌ی مصنوعی با اجرای مدل \\
\lr{ROC} / \lr{zROC curve} & منحنی \lr{ROC} / \lr{zROC} & نرخ \lr{hit}
در برابر نرخ \lr{false alarm} با تغییر ملاک \\
\lr{fit} / \lr{fitting} & برازش & تنظیم پارامترها تا مدل با داده جور
شود \\
\lr{criterion (c)} & ملاک & آستانه‌ی تصمیم؛ نشان‌دهنده‌ی سوگیری پاسخ \\
\lr{formal model} / \lr{formalize} & مدل فرمال / فرمال کردن & نوشتن
نظریه به شکل معادله یا الگوریتم \\
\lr{trial} & کوشش \lr{(trial)} & یک بار انجام تکلیف توسط شرکت‌کننده \\
\lr{false alarm} & هشدار کاذب \lr{(false alarm)} & پاسخ «بله» وقتی علامت
وجود ندارد \\
\lr{estimate} / \lr{estimator} & برآورد / برآوردگر & مقدار به‌دست‌آمده
برای پارامتر / روش به دست آوردن آن \\
\lr{miss} / \lr{correct rejection} & از دست دادن / رد درست & پاسخ «نه»
با وجود علامت / پاسخ «نه» بدون علامت \\
\lr{liberal} / \lr{conservative} & آسان‌گیر / محافظه‌کار & تمایل به گفتن
«بله» / «نه» \\
\lr{free parameter} & پارامتر آزاد & پارامتری که از روی داده برآورد
می‌شود \\
\lr{verbal theory} & نظریه‌ی کلامی & نظریه‌ای که فقط با کلمات بیان شده و
پیش‌بینی دقیق نمی‌دهد \\
\end{longtable}
}

\begin{center}\rule{0.5\linewidth}{0.5pt}\end{center}

\subsection{اهداف یادگیری}

در پایان این هفته، دانشجویان می‌توانند:

\begin{enumerate}
\def\labelenumi{\arabic{enumi}.}
\tightlist
\item
  یک مدل شناختی را به اجزایش \textbf{تجزیه کنند}: طرح آزمایش، پارامترهای
  آزاد و ثابت، هسته‌ی قطعی، مدل نویز (داده)، و فرض‌های پیوند \lr{(linking
  assumptions).}
\item
  \(d'\) و \(c\) را از یک جدول 2×2 شمارش‌ها \textbf{حساب کنند}، از جمله
  با تصحیح برای نرخ \lr{hit} یا \lr{false alarm} برابر 0 یا 1، و
  \textbf{توضیح دهند} چرا \(d'\) حساسیت را از سوگیری جدا می‌کند.
\item
  مدل‌های فراموشی نمایی و توانی را با مدل نویز دوجمله‌ای \textbf{پیاده‌سازی
  کنند} و به‌طور کیفی \textbf{پیش‌بینی کنند} که هر پارامتر منحنی‌ها را
  چگونه تغییر می‌دهد.
\item
  با شبیه‌سازی و نمودارهای تشخیصی \lr{(ROC/zROC}، نمودارهای نیمه‌لگاریتمی
  و لگاریتمی\lr{--}لگاریتمی) \textbf{ارزیابی کنند} که آیا ردپای کیفی یک مدل
  در داده‌ها دیده می‌شود.
\item
  پارامترهای برازش‌شده را به زبان روان‌شناختی \textbf{تفسیر کنند} و
  محدودیت‌های روش «برازش خط راست روی داده‌ی تبدیل‌شده» را \textbf{شناسایی
  کنند}.
\end{enumerate}

\begin{center}\rule{0.5\linewidth}{0.5pt}\end{center}

\subsection{مقدمه}

\textbf{دو رادیولوژیست.} رادیولوژیست \lr{R1} در 93\% از تصویرهایی که
واقعاً تومور دارند تومور گزارش می‌کند، ولی در 31\% از تصویرهای سالم هم.
رادیولوژیست \lr{R2} فقط 69\% از تومورها را گزارش می‌کند، ولی فقط در 7\%
از تصویرهای سالم اشتباه می‌کند. درصد پاسخ درست هر دو یکسان است (81\%).
کدام تشخیص‌دهنده‌ی بهتری است؟

درصد پاسخ درست نمی‌تواند به این سؤال جواب دهد، چون دو چیز متفاوت را قاطی
می‌کند: اینکه فرد \textbf{چقدر خوب} تومور را از بافت سالم تشخیص می‌دهد
(\emph{حساسیت})، و اینکه \textbf{چقدر آمادگی} دارد بگوید «تومور هست»
(\emph{سوگیری}). نظریه‌ی تشخیص علامت یک مدل کوچک دوپارامتری است که این دو
را از هم جدا می‌کند. همان‌طور که مثال حل‌شده نشان می‌دهد، \lr{R1} و \lr{R2}
در واقع حساسیت \textbf{یکسانی} دارند؛ فقط ملاک‌شان را در جای متفاوتی
گذاشته‌اند. این اولین ثمره‌ی مدل‌سازی است: پارامترهای مدل می‌توانند سازه‌های
پنهانی را \textbf{اندازه بگیرند} که نمره‌های خام آن‌ها را قاطی می‌کنند.

\textbf{منحنی ابینگهاوس.} از زمان آزمایش‌هایی که ابینگهاوس در دهه‌ی 1880
روی خودش انجام داد می‌دانیم فراموشی اول سریع است و بعد کند می‌شود. «اول
سریع، بعد کند» توصیفی کلامی است که هم تابع \textbf{نمایی} و هم تابع
\textbf{توانی} با آن جور درمی‌آیند. ولی این دو معناهای متفاوتی دارند:
نمایی یعنی یک خاطره در هر روز \emph{نسبت ثابتی} از قدرتش را از دست
می‌دهد، در هر سنی؛ توانی یعنی خاطره‌های قدیمی‌تر \emph{مقاوم‌تر} می‌شوند.
انتخاب بین این دو سؤالی درباره‌ی ماهیت حافظه است و به یک مدل فرمال نیاز
دارد.

\paragraph{این موضوع کجا به کار می‌رود؟}

\begin{itemize}
\tightlist
\item
  \textbf{تشخیص پزشکی.} رادیولوژیست‌ها، پاتولوژیست‌ها و برنامه‌های غربالگری
  با \lr{SDT} و منحنی \lr{ROC} ارزیابی می‌شوند تا پزشکی که «بهتر می‌بیند»
  از پزشکی که فقط بیشتر می‌گوید «غیرطبیعی است» جدا شود.
\item
  \textbf{پژوهش حافظه.} در آزمون بازشناسی، کلمه‌های مطالعه‌شده «علامت» و
  کلمه‌های جدید «نویز» هستند. \(d'\) قدرت حافظه را جدا از تمایل به گفتن
  «قبلاً دیده‌ام» اندازه می‌گیرد.
\item
  \textbf{شناسایی توسط شاهد عینی.} روش‌های مختلف صف شناسایی
  \lr{(line-up)} در پلیس با تحلیل \lr{ROC} مقایسه می‌شوند: بعضی روش‌ها دقت
  را بیشتر نمی‌کنند و فقط شاهد را محتاط‌تر می‌کنند.
\item
  \textbf{آزمون‌های ادراک و شنوایی.} آستانه‌ی دیدن یا شنیدن محرک‌های ضعیف
  با \lr{SDT} اندازه‌گیری می‌شود تا یک شنونده‌ی محتاط با کسی که کم‌شنواست
  اشتباه گرفته نشود.
\item
  \textbf{یادگیری ماشین.} هر طبقه‌بند دودویی (فیلتر \lr{spam}، آشکارساز
  تقلب، آشکارساز متن هوش مصنوعی) یک نرخ \lr{hit} و یک نرخ \lr{false
  alarm} دارد و با منحنی \lr{ROC} و مساحت زیر آن \lr{(AUC)} ارزیابی
  می‌شود. این همان نظریه است.
\item
  \textbf{منحنی‌های فراموشی} برای زمان‌بندی یادگیری استفاده می‌شوند:
  اپلیکیشن‌های تکرار فاصله‌دار (مثلاً برای لغت) از روی مدلی از سرعت فراموشی
  تصمیم می‌گیرند کِی یک آیتم را دوباره نشان دهند.
\end{itemize}

\begin{center}\rule{0.5\linewidth}{0.5pt}\end{center}

\subsection{جزوه‌ی درس}

\paragraph{به زبان ساده}

این هفته یک ایده‌ی اصلی دارد: هر مدل \textbf{اجزایی} دارد و هر جزء یک فرض
جداست. دو مدل کوچک را می‌بینیم. \textbf{نظریه‌ی تشخیص علامت \lr{(SDT)}}
تصمیم‌های بله/نه در شرایط عدم قطعیت را توضیح می‌دهد. هر \lr{trial} به
مشاهده‌گر یک «حس وجود علامت» پرنویز می‌دهد. اگر این حس از یک آستانه‌ی شخصی
قوی‌تر باشد، مشاهده‌گر می‌گوید «بله». دو عدد مشاهده‌گر را توصیف می‌کنند: حس
«علامت هست» و «علامت نیست» به‌طور میانگین چقدر از هم فاصله دارند
(\textbf{حساسیت}، \(d'\))، و آستانه کجاست (\textbf{ملاک}، \(c\)). دو نفر
می‌توانند درصد پاسخ درست یکسانی داشته باشند به دلایل کاملاً متفاوت، و
\lr{SDT} نشان می‌دهد کدام دلیل. \textbf{منحنی فراموشی} نشان می‌دهد یادآوری
با گذشت زمان چطور کم می‌شود. دو فرمول، نمایی و توانی، هر دو شکل «اول
سریع، بعد آهسته» دارند، ولی درباره‌ی حافظه حرف‌های متفاوتی می‌زنند. انتخاب
بین آن‌ها از جنس سؤال‌هایی است که بقیه‌ی درس یاد می‌دهد با داده جواب دهیم.

\paragraph{واژه‌های آماری به زبان ساده}

\begin{itemize}
\tightlist
\item
  \textbf{توزیع \lr{(distribution):}} توصیف اینکه یک کمیت تصادفی چه
  مقدارهایی می‌گیرد و هرکدام چقدر رایج است. هیستوگرام تعداد زیادی
  اندازه‌گیری تقریبی از آن است.
\item
  \textbf{توزیع نرمال (گاوسی) \(\mathcal{N}(\mu, \sigma)\):} منحنی
  زنگوله‌ای با میانگین \(\mu\) (مرکز) و انحراف معیار \(\sigma\) (پهنا).
  حدود 68\% مقدارها در \(\mu \pm \sigma\) و 95\% در
  \(\mu \pm 1.96\sigma\) قرار می‌گیرند.
\item
  \textbf{نرمال استاندارد:} توزیع نرمال با \(\mu = 0\) و \(\sigma = 1\).
  مقداری که روی این مقیاس سنجیده شود \textbf{\lr{z-score}} نام دارد:
  «چند انحراف معیار بالاتر از میانگین».
\item
  \textbf{\(\Phi(x)\) (تابع توزیع تجمعی یا \lr{CDF):}} احتمال اینکه یک
  مقدار نرمال استاندارد \emph{کمتر} از \(x\) باشد. \(\Phi(0) = 0.5\)،
  \(\Phi(1) \approx 0.84\)، \(\Phi(1.96) \approx 0.975\). در
  \lr{Python:} \lr{\nolinkurl{norm.cdf(x)}}.
\item
  \textbf{\(\Phi^{-1}(p)\) (معکوس \lr{CDF}، «تبدیل \lr{z»):}} سؤال
  برعکس: نسبت \(p\) از مقدارها زیر چه عددی هستند؟
  \(\Phi^{-1}(0.84) \approx 1\). در \lr{Python:}
  \lr{\nolinkurl{norm.ppf(p)}}.
\item
  \textbf{توزیع دوجمله‌ای \(\mathrm{Binomial}(n, p)\):} تعداد موفقیت‌ها در
  \(n\) تلاش مستقل که هرکدام با احتمال \(p\) موفق است (مثلاً تعداد
  کلمه‌های به‌یادآمده از 50). میانگین آن \(np\) و واریانس آن \(np(1-p)\)
  است.
\item
  \textbf{نسبت در برابر احتمال:} نسبت چیزی است که \emph{مشاهده} می‌کنیم
  (30 از 50 = 0.60)؛ احتمال مقدار \emph{زیربنایی} مدل است که نسبت آن را
  برآورد می‌کند.
\end{itemize}

\textbf{زمان‌بندی پیشنهادی (60 دقیقه):} کالبدشکافی مدل 15 · \lr{SDT} 20 ·
منحنی‌های فراموشی 15 · شبیه‌سازی به‌عنوان ابزار فکر کردن + پیش‌نمایش 10.

\subsubsection{1. کالبدشکافی یک مدل}

در هفته‌ی 1 مدل را یک ماشین تولید داده نوشتیم،
\(\mathbf{y} \sim p(\mathbf{y}\mid\boldsymbol\theta,\allowbreak \mathcal{M})\). در
عمل تقریباً هر مدلی که در این درس می‌بینید را می‌شود به پنج تکه تقسیم کرد:

{\def\LTcaptype{none} % do not increment counter
\begin{longtable}[]{@{}
  >{\raggedright\arraybackslash}p{(\linewidth - 6\tabcolsep) * \real{0.2500}}
  >{\raggedright\arraybackslash}p{(\linewidth - 6\tabcolsep) * \real{0.2500}}
  >{\raggedright\arraybackslash}p{(\linewidth - 6\tabcolsep) * \real{0.2500}}
  >{\raggedright\arraybackslash}p{(\linewidth - 6\tabcolsep) * \real{0.2500}}@{}}
\toprule\noalign{}
\begin{minipage}[b]{\linewidth}\raggedright
جزء
\end{minipage} & \begin{minipage}[b]{\linewidth}\raggedright
سؤال
\end{minipage} & \begin{minipage}[b]{\linewidth}\raggedright
\lr{SDT}
\end{minipage} & \begin{minipage}[b]{\linewidth}\raggedright
فراموشی (نمایی)
\end{minipage} \\
\midrule\noalign{}
\endhead
\bottomrule\noalign{}
\endlastfoot
\textbf{1. طرح آزمایش} \(x_i\) & چه چیزی دست‌کاری و اندازه‌گیری می‌شود؟ &
نوع \lr{trial} (علامت/نویز)؛ پاسخ بله/نه & فاصله‌ی نگهداری \(t_i\)؛ تعداد
آیتم‌های یادآوری‌شده از \(n\) آیتم \\
\textbf{2. پارامترها} \(\boldsymbol\theta\) & چه چیزی آزاد است؟ چه چیزی
ثابت؟ & آزاد: \(d'\)، \(c\). ثابت: هر دو انحراف معیار = 1، میانگین نویز
= 0 & آزاد: \(A\)، \(\lambda\). ثابت: \(n\) \\
\textbf{3. هسته‌ی قطعی} \(f(x;\boldsymbol\theta)\) & مدل به‌طور میانگین چه
پیش‌بینی می‌کند؟ & \(H = \Phi(d'/2 - c)\)، \(F = \Phi(-d'/2 - c)\) &
\(f(t) = A e^{-\lambda t}\) \\
\textbf{4. مدل نویز} & مشاهدات چگونه حول پیش‌بینی پراکنده می‌شوند؟ &
\lr{hit}ها \(\sim \mathrm{Binomial}(n_s, H)\)؛ \lr{false alarm}ها
\(\sim \mathrm{Binomial}(n_n, F)\) & یادآوری‌شده‌ها
\(\sim \mathrm{Binomial}(n, f(t))\) \\
\textbf{5. فرض‌های پیوند} & کمیت‌های مدل چگونه به رفتار نگاشت می‌شوند؟ &
اگر شواهد درونی از ملاک بیشتر باشد «بله» & احتمال یادآوری برابر قدرت
حافظه \(f(t)\) است \\
\end{longtable}
}

چند نکته درباره‌ی هر جزء:

\begin{itemize}
\tightlist
\item
  \textbf{پارامترهای آزاد در برابر ثابت.} پارامترهای ثابت «بی‌اهمیت»
  نیستند؛ آن‌ها \emph{فرض} هستند. در هفته‌ی 1، \(t_0 = 300\) \lr{ms} را
  ثابت گرفتیم. در \lr{SDT}، ثابت گرفتن میانگین نویز روی 0 و هر دو انحراف
  معیار روی 1 است که اصلاً برآورد \(d'\) و \(c\) را ممکن می‌کند: جابه‌جا
  کردن یا کشیدن کل محور هیچ پیش‌بینی‌ای را عوض نمی‌کند. به چنین ثابت‌هایی
  \textbf{قیدهای مقیاس} \lr{(scaling constraints)} می‌گویند.
\item
  \textbf{دامنه‌ی پارامترها.} هر پارامتر دامنه‌ی معناداری دارد:
  \(0 \le A \le 1\)، \(\lambda > 0\)، و برای مشاهده‌گری که برعکس پاسخ
  نمی‌دهد \(d' \ge 0\). رعایت این دامنه‌ها در برازش مهم است (هفته‌ی 3).
\item
  \textbf{مدل نویز جزئی از مدل است.} دو پژوهشگر با \(f(t)\) یکسان ولی
  مدل نویز متفاوت (دوجمله‌ای در برابر گاوسی) مدل‌های \emph{متفاوتی} دارند
  و برآوردهای پارامتر متفاوتی می‌گیرند (هفته‌ی 4). توزیع دوجمله‌ای برای
  شمارش‌ها طبیعی است: واریانسش \(f(1-f)/n\) در \(f = 0.5\) بیشینه است و
  نزدیک 0 و 1 کم می‌شود.
\end{itemize}

\subsubsection{\texorpdfstring{2. نظریه‌ی تشخیص علامت
\lr{(SDT)}}{2. نظریه‌ی تشخیص علامت (SDT)}}

\textbf{تکلیف آزمایشی.} در هر \lr{trial} محرک یا \emph{علامت} است (مثلاً
یک تومور، یک کلمه‌ی مطالعه‌شده، یک صدای ضعیف) یا \emph{نویز}. مشاهده‌گر
پاسخ «بله» (علامت) یا «نه» می‌دهد. نتیجه‌ها یک جدول 2×2 می‌سازند:

{\def\LTcaptype{none} % do not increment counter
\begin{longtable}[]{@{}lll@{}}
\toprule\noalign{}
& پاسخ «بله» & پاسخ «نه» \\
\midrule\noalign{}
\endhead
\bottomrule\noalign{}
\endlastfoot
\textbf{\lr{trial} علامت} & \lr{hit} & از دست دادن \lr{(miss)} \\
\textbf{\lr{trial} نویز} & \lr{false alarm (FA)} & رد درست \lr{(correct
rejection}، \lr{CR)} \\
\end{longtable}
}

با نرخ \lr{hit} \(H = \frac{\text{\lr{hits}}}{\text{\lr{hits}} + \text{\lr{misses}}}\)
و نرخ \lr{false alarm}
\(F = \frac{\text{\lr{FAs}}}{\text{\lr{FAs}} + \text{\lr{CRs}}}\).

\textbf{مدل (شهودی).} هر \lr{trial} یک مقدار \textbf{شواهد} درونی تولید
می‌کند. این مقدار پرنویز است: حتی \lr{trial}های نویز خالص هم گاهی شبیه
علامت به نظر می‌رسند. \lr{trial}های علامت به‌طور میانگین شواهد بیشتری
تولید می‌کنند. مشاهده‌گر هر وقت شواهد از یک \textbf{ملاک} \lr{(criterion)}
بیشتر باشد می‌گوید «بله».

\textbf{مدل (فرمال)، \lr{SDT} گاوسی با واریانس برابر:}

\[
\text{\lr{noise trials: }} X \sim \mathcal{N}(0, 1), \qquad
\text{\lr{signal trials: }} X \sim \mathcal{N}(d', 1), \qquad
\text{\lr{respond "yes" iff }} X > k.
\]

\begin{itemize}
\tightlist
\item
  \(d'\) \textbf{حساسیت} است: فاصله‌ی دو توزیع بر حسب واحد انحراف معیار.
\item
  \(k\) محل مطلق ملاک است. راحت‌تر است ملاک را نسبت به نقطه‌ی میانی دو
  توزیع بسنجیم: \(c = k - d'/2\). با \(c = 0\) مشاهده‌گر بی‌سوگیری است،
  \(c < 0\) یعنی \textbf{\lr{liberal}} (راحت «بله» می‌گوید)، و \(c > 0\)
  یعنی \textbf{محافظه‌کار \lr{(conservative)}}.
\end{itemize}

\textbf{پیش‌بینی‌ها.} با استفاده از \lr{CDF} نرمال استاندارد \(\Phi\):

\[
H = P(X > k \mid \text{\lr{signal}}) = \Phi\!\left(\tfrac{d'}{2} - c\right), \qquad
F = P(X > k \mid \text{\lr{noise}}) = \Phi\!\left(-\tfrac{d'}{2} - c\right).
\]

\textbf{معکوس کردن مدل.} با اعمال \(\Phi^{-1}\) (که اغلب \(z\) نوشته
می‌شود) روی هر دو معادله و حل آن‌ها، فرمول‌های آشنا به دست می‌آیند:

\[
d' = \Phi^{-1}(H) - \Phi^{-1}(F), \qquad
c = -\tfrac{1}{2}\left[\Phi^{-1}(H) + \Phi^{-1}(F)\right].
\]

دقت کنید چه اتفاقی افتاد: این اولین \textbf{برآورد پارامتر} ماست. دو
پارامتر آزاد و دو داده (\(H\)، \(F\)) داریم، پس مدل را می‌شود دقیقاً حل
کرد. وقتی تعداد داده‌ها از تعداد پارامترها بیشتر باشد، معمولاً جواب دقیق
وجود ندارد و به روش‌های برازش هفته‌های 3\lr{--}4 نیاز داریم.

\textbf{مشکل نرخ‌های حدی.} اگر \(H = 1\) یا \(F = 0\) باشد، \(\Phi^{-1}\)
برابر \(\pm\infty\) و \(d'\) بی‌نهایت می‌شود. با تعداد \lr{trial} کم این
اتفاق زیاد می‌افتد. راه‌حل استاندارد \textbf{تصحیح \lr{log-linear}} است:
پیش از محاسبه‌ی نرخ‌ها به هر خانه‌ی جدول 0.5 اضافه کنید (یعنی 0.5 به
\lr{hit}ها و \lr{false alarm}ها و 1 به تعداد \lr{trial}های هر نوع). در
\lr{lab} به‌طور پیش‌فرض از آن استفاده می‌کنیم.

\textbf{منحنی \lr{ROC:} نشانه‌ی اصلی \lr{SDT.}} \(d'\) را ثابت نگه دارید
و بگذارید \(c\) تغییر کند. جفت \((F, H)\) یک منحنی رسم می‌کند به نام
\textbf{منحنی مشخصه‌ی عملکرد گیرنده \lr{(ROC)}}. همه‌ی نقاط یک \lr{ROC}
حساسیت یکسان ولی سوگیری متفاوت دارند. رسم \(\Phi^{-1}(H)\) در برابر
\(\Phi^{-1}(F)\) منحنی \textbf{\lr{zROC}} را می‌دهد که در \lr{SDT} با
واریانس برابر یک خط راست است:

\[
\Phi^{-1}(H) = d' + 1 \cdot \Phi^{-1}(F).
\]

پس مدل یک پیش‌بینی \emph{قابل آزمون} دارد: \textbf{\lr{zROC} خطی با شیب 1
است}. \lr{ROC} را می‌شود با درجه‌بندی اطمینان \lr{(confidence ratings)}
اندازه گرفت (هر سطح اطمینان مثل یک ملاک جداگانه عمل می‌کند). در حافظه‌ی
بازشناسی، شیب \lr{zROC} معمولاً \emph{کمتر} از 1 است و همین انگیزه‌ی مدل
با واریانس نابرابر در ادامه است.

\begin{quote}
\textbf{پیشرفته: \lr{SDT} با واریانس نابرابر.} اگر در \lr{trial}های
علامت \(X \sim \mathcal{N}(\mu_s,\allowbreak \sigma_s)\) باشد، \lr{zROC} هنوز خطی
است ولی شیبش \(1/\sigma_s\) است. در این حالت \(d'\) که از یک جفت
\((H, F)\) حساب می‌شود \emph{با ملاک تغییر می‌کند}، پس دیگر یک معیار خالص
حساسیت نیست. یک شاخص جایگزین رایج
\(d_a = \mu_s \sqrt{2/(1+\sigma_s^2)}\) است. برآورد \(\sigma_s\) به بیش
از یک ملاک، یعنی داده‌ی درجه‌بندی، نیاز دارد. این اثر را مستقیماً در بخش
\lr{B lab} می‌بینید.
\end{quote}

\paragraph{خواندن ریاضی}

\textbf{دو منحنی زنگوله‌ای.} دو منحنی نرمال را روی یک محور بکشید: منحنی
«نویز» با مرکز 0 و منحنی «علامت» با مرکز \(d'\). پهنای هر دو 1 است. یک
خط عمودی در محل ملاک \(k\) بکشید. هر چیزی سمت راست خط یعنی «بله».

\begin{itemize}
\tightlist
\item
  \(H\) = مساحت منحنی \textbf{علامت} در سمت راست خط. \(F\) = مساحت منحنی
  \textbf{نویز} در سمت راست خط.
\item
  چرا \(H = \Phi(d'/2 - c)\)؟ ملاک در \(k = d'/2 + c\) است (نقطه‌ی وسط
  به‌علاوه‌ی \(c\)). مرکز منحنی علامت \(d'\) است، پس خط به اندازه‌ی
  \(k - d' = c - d'/2\) انحراف معیار از مرکز آن فاصله دارد. مساحت سمت
  راست این نقطه \(1 - \Phi(c - d'/2) = \Phi(d'/2 - c)\) است، چون منحنی
  نرمال متقارن است.
\item
  \textbf{مثال عددی.} \(d' = 2\) و \(c = 0\): \(H = \Phi(1) = 0.84\) و
  \(F = \Phi(-1) = 0.16\). حالا مشاهده‌گر را آسان‌گیر \lr{(liberal)} کنید،
  \(c = -0.5\): \(H = \Phi(1.5) = 0.93\) و \(F = \Phi(-0.5) = 0.31\). هر
  دو نرخ بالا رفتند و حساسیت تغییر نکرد. این همان رادیولوژیست \lr{R1} در
  مثال حل‌شده است.
\end{itemize}

\textbf{معکوس کردن.} \(d' = \Phi^{-1}(H) - \Phi^{-1}(F)\): هر نرخ را به
\lr{z-score} تبدیل کنید و از هم کم کنید. برای \lr{R1:}
\(1.476 - (-0.496) = 1.97\). \(c\) منفیِ میانگین دو \lr{z-score} است: اگر
هر دو نرخ بالا باشند (هر دو \lr{z-score} مثبت)، \(c\) منفی است، یعنی
آسان‌گیر.

\textbf{خط \lr{zROC.}} \(\Phi^{-1}(H) = d' + \Phi^{-1}(F)\) همان معادله‌ی
قبلی است که مرتب شده. در واحد \lr{z}، جابه‌جا کردن ملاک هر دو
\lr{z-score} را به یک اندازه تغییر می‌دهد، پس نقطه‌ها روی خطی با شیب 1 و
عرض از مبدأ \(d'\) قرار می‌گیرند.

\paragraph{برای یک مهندس}

\lr{SDT} \emph{همان} نظریه‌ی آشکارسازی رادار و مخابرات است که در دهه‌ی
1950 وارد روان‌شناسی شد.

{\def\LTcaptype{none} % do not increment counter
\begin{longtable}[]{@{}
  >{\raggedright\arraybackslash}p{(\linewidth - 2\tabcolsep) * \real{0.5000}}
  >{\raggedright\arraybackslash}p{(\linewidth - 2\tabcolsep) * \real{0.5000}}@{}}
\toprule\noalign{}
\begin{minipage}[b]{\linewidth}\raggedright
روان‌شناسی
\end{minipage} & \begin{minipage}[b]{\linewidth}\raggedright
مخابرات / رادار
\end{minipage} \\
\midrule\noalign{}
\endhead
\bottomrule\noalign{}
\endlastfoot
\lr{trial} نویز / \lr{trial} علامت & فرضیه‌ی \(H_0\) (فقط نویز) / \(H_1\)
(سیگنال + نویز) \\
شواهد درونی \(X\) & آماره‌ی آزمون، مثلاً خروجی فیلتر منطبق \lr{(matched
filter)} \\
ملاک \(k\) (یا \(c\)) & آستانه‌ی آشکارسازی \\
نرخ \lr{hit} یعنی \(H\) & احتمال آشکارسازی \(P_D\) \\
نرخ \lr{false alarm} یعنی \(F\) & احتمال هشدار کاذب \(P_{FA}\) \\
حساسیت \(d'\) & انحراف \lr{(deflection)} \((\mu_1 - \mu_0)/\sigma\)؛
\(d'^2\) نقش \lr{SNR} را دارد \\
منحنی \lr{ROC} & \lr{receiver operating characteristic} (همان نام، همان
نمودار) \\
مشاهده‌گر آسان‌گیر / محافظه‌کار & آستانه‌ی پایین / بالا (\(P_{FA}\) بالا /
پایین) \\
\end{longtable}
}

تفاوت در تأکید: مهندس معمولاً آستانه را \emph{انتخاب} می‌کند (مثلاً
\lr{Neyman\lr{--}Pearson:} \(P_{FA}\) را ثابت کن و \(P_D\) را بیشینه کن).
روان‌شناس \(P_D\) و \(P_{FA}\) یک آدم را \emph{مشاهده} می‌کند و هم
\lr{SNR} (\(d'\)) و هم آستانه‌ای (\(c\)) را که آن فرد به کار برده
\emph{استنباط} می‌کند. ملاک \(c\) از نقطه‌ی وسط دو میانگین سنجیده می‌شود،
پس \(c = 0\) یعنی «بدون سوگیری». آستانه‌ی \textbf{بهینه} دقیقاً مثل
آشکارسازی بیزی به احتمال‌های پیشین و هزینه‌ها بستگی دارد، و آدم‌ها وقتی
پاداش‌ها تغییر می‌کنند معمولاً ملاکشان را در جهت درست جابه‌جا می‌کنند (هفته‌ی
4).

\subsubsection{3. منحنی‌های فراموشی}

\textbf{داده.} در هر فاصله‌ی نگهداری \(t_i\) (مثلاً به روز)، شرکت‌کننده با
\(n\) آیتم آزمون می‌شود و \(y_i\) تا را به یاد می‌آورد.

\textbf{دو هسته‌ی قطعی رقیب:}

{\def\LTcaptype{none} % do not increment counter
\begin{longtable}[]{@{}
  >{\raggedright\arraybackslash}p{(\linewidth - 4\tabcolsep) * \real{0.3333}}
  >{\raggedright\arraybackslash}p{(\linewidth - 4\tabcolsep) * \real{0.3333}}
  >{\raggedright\arraybackslash}p{(\linewidth - 4\tabcolsep) * \real{0.3333}}@{}}
\toprule\noalign{}
\begin{minipage}[b]{\linewidth}\raggedright
\end{minipage} & \begin{minipage}[b]{\linewidth}\raggedright
نمایی
\end{minipage} & \begin{minipage}[b]{\linewidth}\raggedright
توانی
\end{minipage} \\
\midrule\noalign{}
\endhead
\bottomrule\noalign{}
\endlastfoot
معادله & \(f(t) = A\,e^{-\lambda t}\) & \(f(t) = A\,(1 + t)^{-d}\) \\
پارامترها & \(A\): سطح اولیه؛ \(\lambda\): نرخ زوال & \(A\): سطح اولیه؛
\(d\): نمای زوال \\
نمودار خطی‌کننده & \(\ln f\) در برابر \(t\) خطی راست با شیب \(-\lambda\)
است & \(\ln f\) در برابر \(\ln(1+t)\) خطی راست با شیب \(-d\) است \\
نرخ نسبی فراموشی \(-f'(t)/f(t)\) & \(\lambda\) (ثابت) & \(d/(1+t)\) (با
گذشت زمان کاهشی) \\
خوانش روان‌شناختی & خاطره در هر واحد زمان \emph{نسبت} یکسانی از قدرتش را
از دست می‌دهد، هر سنی که داشته باشد & خاطره‌های قدیمی‌تر کندتر قدرت از دست
می‌دهند؛ مقاوم‌تر می‌شوند \\
\end{longtable}
}

\textbf{نرخ نسبی فراموشی} (که در تحلیل بقا به آن \emph{\lr{hazard}}
می‌گویند) تفاوت کلیدی در تفسیر است. به همین دلیل انتخاب بین این دو تابع
یک سؤال نظری است، نه یک جزئیات منحنی‌برازی.

\textbf{مدل نویز.} تعداد یادآوری‌ها دوجمله‌ای است:

\[
y_i \sim \mathrm{Binomial}\big(n, f(t_i; \boldsymbol\theta)\big).
\]

با \(n = 50\) آیتم و \(f = 0.5\)، انحراف معیار نسبت مشاهده‌شده
\(\sqrt{0.5 \cdot 0.5 / 50} \approx 0.07\) است. این در مقایسه با تفاوت
مدل‌ها در فاصله‌های کوتاه بزرگ است؛ اولین نشانه‌ی اینکه \textbf{طرح آزمایش}
(کدام فاصله‌ها، چند آیتم) تعیین می‌کند داده‌ها می‌توانند بین مدل‌ها تمایز
بگذارند یا نه.

\begin{quote}
\textbf{پیشرفته: «+1» در مدل توانی بی‌خطر نیست.} \(A\,t^{-d}\) در
\(t = 0\) بی‌نهایت است، پس یک جابه‌جایی اضافه می‌شود. ولی \((1 + t)\) به
\textbf{واحد} \(t\) بستگی دارد: اندازه‌گیری به ساعت به‌جای روز، شکل منحنی
را عوض می‌کند، نه فقط مقیاس محور را. به همین دلیل بعضی نویسندگان
\(A\,(1 + bt)^{-d}\) را با پارامتر مقیاس آزاد \(b\) می‌نویسند؛ این یک
پارامتر سوم اضافه می‌کند و سؤال‌های مقایسه‌ی مدل هفته‌های 6\lr{--}7 را پیش
می‌آورد. در طول این درس واحد را ثابت (روز) نگه دارید.
\end{quote}

\paragraph{خواندن ریاضی}

\begin{itemize}
\tightlist
\item
  \(f(t) = A e^{-\lambda t}\): \(A\) میزان یادآوری درست بعد از مطالعه
  (\(t = 0\)) است و \(\lambda\) می‌گوید با چه سرعتی کم می‌شود. با
  \(\lambda = 0.1\) در روز، یادآوری هر روز در \(e^{-0.1} \approx 0.905\)
  ضرب می‌شود: هر روز حدود 9.5\% از آنچه مانده از دست می‌رود، چه روز اول و
  چه روز صدم.
\item
  \(f(t) = A(1+t)^{-d}\): رفتن از روز 0 به روز 1 یادآوری را در
  \(2^{-d}\) ضرب می‌کند؛ رفتن از روز 9 به روز 10 فقط در \((10/11)^{d}\)
  که نزدیک 1 است. خاطره‌های قدیمی دیگر تقریباً تغییر نمی‌کنند.
\item
  \textbf{نرخ نسبی فراموشی} \(-f'(t)/f(t)\): «چه کسری از حافظه‌ی
  باقی‌مانده \emph{همین الان} در هر روز از دست می‌رود». برای تابع نمایی
  ثابت است (\(\lambda\)) و برای تابع توانی کم می‌شود (\(d/(1+t)\)). همین
  یک عدد، تفاوت روان‌شناختی دو مدل است.
\item
  \(y_i \sim \mathrm{Binomial}(n,\allowbreak f(t_i))\): در تأخیر \(t_i\)، شرکت‌کننده
  \(y_i\) آیتم از \(n\) آیتم را به یاد می‌آورد و هر آیتم با احتمال
  \(f(t_i)\) به یاد آورده می‌شود. نسبت مشاهده‌شده‌ی \(y_i/n\) با انحراف
  معیار \(\sqrt{f(1-f)/n}\) حول \(f(t_i)\) پراکنده است؛ برای \(n = 50\)
  حدود 0.07.
\end{itemize}

\paragraph{برای یک مهندس}

\begin{itemize}
\tightlist
\item
  منحنی \textbf{نمایی} همان پاسخ پله‌ی یک سیستم مرتبه‌ی اول است (دشارژ یک
  مدار \lr{RC):} یک ثابت زمانی \(\tau = 1/\lambda\). ثابت بودن نرخ نسبی
  فراموشی همان ویژگی «بی‌حافظگی» است که از طول عمر نمایی در نظریه‌ی قابلیت
  اطمینان می‌شناسید.
\item
  منحنی \textbf{توانی} یک ثابت زمانی واحد ندارد. تصویر مفید: جمع تعداد
  زیادی منحنی نمایی با ثابت‌های زمانی \emph{متفاوت} (مؤلفه‌های سریع و کند)
  شبیه یک قانون توانی می‌شود. هفته‌ی 5 این را دقیقاً نشان می‌دهد و این یکی
  از توضیح‌های ممکن برای فراموشی توانی است.
\item
  روی نمودار \textbf{لگاریتمی\lr{--}لگاریتمی} قانون توانی یک خط راست است؛ روی
  نمودار \textbf{نیمه‌لگاریتمی} تابع نمایی خط راست است. مهندس‌ها از همین
  ترفند برای تشخیص مرتبه‌ی سیستم و نویز \(1/f\) استفاده می‌کنند.
\item
  مدل نویز دوجمله‌ای همان آمار شمارش یک آشکارساز است: مثل شمردن خطاهای
  بیتی از \(n\) بیت، عدم قطعیت یک نسبت با \(1/\sqrt{n}\) کم می‌شود.
\end{itemize}

\subsubsection{4. شبیه‌سازی به‌عنوان ابزار فکر کردن}

پیش از برازش هر چیزی، شبیه‌سازی کنید. شبیه‌سازی به سؤال‌هایی مثل این‌ها جواب
می‌دهد:

\begin{itemize}
\tightlist
\item
  \textbf{مدل چه چیزی را می‌تواند تولید کند، و چه چیزی را
  \emph{نمی‌تواند}؟} \lr{SDT} با واریانس برابر و \(d' > 0\) هرگز نمی‌تواند
  \(H < F\) تولید کند، و هرگز شیب \lr{zROC} غیر از 1 نمی‌دهد. منحنی‌های
  نمایی و توانی هر دو همیشه نزولی‌اند، پس نمی‌توانند \emph{بازیادآوری}
  \lr{(reminiscence}، بهتر شدن یادآوری با گذشت زمان) را نشان دهند.
  مجموعه‌ی کارهایی که یک مدل \textbf{نمی‌تواند} انجام دهد همان چیزی است که
  آن را قابل آزمون می‌کند.
\item
  \textbf{کدام پارامتر کجا اهمیت دارد؟} \(A\) کل منحنی فراموشی را جابه‌جا
  می‌کند؛ \(\lambda\) یا \(d\) شکل را کنترل می‌کند.
\item
  \textbf{آیا طرح آزمایش می‌تواند بین مدل‌ها تمایز بگذارد؟} اگر دو مدل در
  فاصله‌هایی که می‌خواهیم آزمون کنیم پیش‌بینی‌های تقریباً یکسانی بدهند، هیچ
  روش برازشی نمی‌تواند آن‌ها را از هم جدا کند (تمرین 4 \lr{lab}؛ تکلیف
  \lr{A1).}
\end{itemize}

\subsubsection{5. پیش‌نمایش: برآورد با خطی‌سازی، و چرا کافی نیست}

یک روش سنتی برای برازش منحنی فراموشی این است که از نسبت‌های مشاهده‌شده
لگاریتم بگیریم و با رگرسیون معمولی یک خط راست برازش کنیم. سریع و شهودی
است و در \lr{lab} امتحانش می‌کنیم. سه مشکل دارد که انگیزه‌ی هفته‌های 3\lr{--}4
است:

\begin{enumerate}
\def\labelenumi{\arabic{enumi}.}
\tightlist
\item
  \(\ln(0)\) تعریف‌نشده است، پس فاصله‌ای با یادآوری صفر روش را از کار
  می‌اندازد.
\item
  تبدیل لگاریتمی نویز را تحریف می‌کند: پراکندگی \(\ln \hat p\) برای
  نسبت‌های کوچک خیلی بیشتر است، ولی رگرسیون با همه‌ی نقاط مثل هم رفتار
  می‌کند.
\item
  خط روی مقیاس لگاریتمی برازش می‌شود، ولی ما به خطا روی مقیاس اصلی (نسبت)
  اهمیت می‌دهیم.
\end{enumerate}

جایگزین اصولی این است که مدل را مستقیماً روی داده‌ی تبدیل‌نشده برازش کنیم:
با کمینه کردن مربع خطا (هفته‌ی 3) یا بیشینه کردن درست‌نمایی دوجمله‌ای
(هفته‌ی 4).

\paragraph{برای مدرس}

\begin{itemize}
\tightlist
\item
  \textbf{با معمای رادیولوژیست‌ها شروع کنید} (بخش مقدمه) و پیش از نشان
  دادن هر فرمولی، بگذارید دانشجوها رأی بدهند کدام بهتر است. بیشترشان بر
  اساس درصد پاسخ درست انتخاب می‌کنند و بعد می‌فهمند که دو عدد یکسان است.
\item
  \textbf{رسم کنید، اثبات نکنید.} دو منحنی زنگوله‌ای هم‌پوشان با یک خط
  عمودی متحرک، \(H\)، \(F\)، \(d'\) و \(c\) را بهتر از هر فرمولی توضیح
  می‌دهند. خط را روی تخته به چپ و راست ببرید و بپرسید \(H\) و \(F\) چه
  می‌شوند؛ بعد منحنی‌ها را از هم دور کنید و دوباره بپرسید.
\item
  \textbf{از جدول \(\Phi\) یا \lr{Python} به‌طور زنده استفاده کنید.}
  \(d'\) رادیولوژیست \lr{R1} را جلوی کلاس با \lr{\nolinkurl{norm.ppf}}
  حساب کنید. دانشجوهایی که آمار برایشان تازه است باید ببینند
  \(\Phi^{-1}\) «فقط یک جدول‌خوانی» است.
\item
  \textbf{به چیزی که می‌شناسند وصل کنید:} یک تست کرونا یا بارداری نرخ
  \lr{hit (sensitivity)} و نرخ \lr{false alarm} (یک منهای
  \lr{specificity)} دارد. \lr{«sensitivity»} در پزشکی همان نرخ \lr{hit}
  است، \emph{نه} \(d'\): به این تشابه اسمی اشاره کنید.
\item
  \textbf{سؤال‌های رایج.}

  \begin{itemize}
  \tightlist
  \item
    \emph{«چرا توزیع نرمال فرض می‌کنیم؟»} فرضی راحت و اغلب کافی است؛
    آزمون \lr{zROC} (شیب 1) آن را بررسی می‌کند و مدل واریانس نابرابر آن
    را آزادتر می‌کند.
  \item
    \emph{«آیا ملاک آسان‌گیر بد است؟»} نه؛ بهترین ملاک به هزینه‌ها و
    نرخ‌های پایه بستگی دارد.
  \item
    \emph{«نمایی یا توانی، کدام درست است؟»} در بسیاری از داده‌ها تابع
    توانی بهتر برازش می‌شود، ولی میانگین‌گیری روی افراد می‌تواند شکل توانی
    بسازد (هفته‌ی 5). سؤال را باز نگه دارید.
  \end{itemize}
\end{itemize}

\begin{center}\rule{0.5\linewidth}{0.5pt}\end{center}

\subsection{مثال حل‌شده}

\subsubsection{بخش 1: دو رادیولوژیست (و یک سومی)}

هر رادیولوژیست 100 تصویر با تومور و 100 تصویر سالم را می‌خواند.

{\def\LTcaptype{none} % do not increment counter
\begin{landscape}\small
\begin{longtable}[]{@{}
  >{\raggedright\arraybackslash}p{(\linewidth - 14\tabcolsep) * \real{0.1250}}
  >{\raggedright\arraybackslash}p{(\linewidth - 14\tabcolsep) * \real{0.1250}}
  >{\raggedright\arraybackslash}p{(\linewidth - 14\tabcolsep) * \real{0.1250}}
  >{\raggedright\arraybackslash}p{(\linewidth - 14\tabcolsep) * \real{0.1250}}
  >{\raggedright\arraybackslash}p{(\linewidth - 14\tabcolsep) * \real{0.1250}}
  >{\raggedright\arraybackslash}p{(\linewidth - 14\tabcolsep) * \real{0.1250}}
  >{\raggedright\arraybackslash}p{(\linewidth - 14\tabcolsep) * \real{0.1250}}
  >{\raggedright\arraybackslash}p{(\linewidth - 14\tabcolsep) * \real{0.1250}}@{}}
\toprule\noalign{}
\begin{minipage}[b]{\linewidth}\raggedright
مشاهده‌گر
\end{minipage} & \begin{minipage}[b]{\linewidth}\raggedright
\(H\)
\end{minipage} & \begin{minipage}[b]{\linewidth}\raggedright
\(F\)
\end{minipage} & \begin{minipage}[b]{\linewidth}\raggedright
\(\Phi^{-1}(H)\)
\end{minipage} & \begin{minipage}[b]{\linewidth}\raggedright
\(\Phi^{-1}(F)\)
\end{minipage} & \begin{minipage}[b]{\linewidth}\raggedright
\(d'\)
\end{minipage} & \begin{minipage}[b]{\linewidth}\raggedright
\(c\)
\end{minipage} & \begin{minipage}[b]{\linewidth}\raggedright
درصد درست
\end{minipage} \\
\midrule\noalign{}
\endhead
\bottomrule\noalign{}
\endlastfoot
\lr{R0} (خنثی) & 0.84 & 0.16 & 0.994 & \ensuremath{-}0.994 & 1.99 & 0.00 & 84\% \\
\lr{R1 (liberal)} & 0.93 & 0.31 & 1.476 & \ensuremath{-}0.496 & 1.97 & \ensuremath{-}0.49 &
81\% \\
\lr{R2} (محافظه‌کار) & 0.69 & 0.07 & 0.496 & \ensuremath{-}1.476 & 1.97 & +0.49 &
81\% \\
\end{longtable}
\end{landscape}
}

\textbf{قدم‌به‌قدم برای \lr{R1:}} 1. \(\Phi^{-1}(0.93) = 1.476\) و
\(\Phi^{-1}(0.31) = -0.496\) (\lr{\nolinkurl{scipy.stats.norm.ppf}}). 2.
\(d' = 1.476 - (-0.496) = 1.972\). 3.
\(c = -\tfrac{1}{2}(1.476 + (-0.496)) = -0.490\)، که منفی است، پس \lr{R1
liberal} است. 4. بررسی با جای‌گذاری:
\(H = \Phi(1.972/2 + 0.490) = \Phi(1.476) = 0.93\). \ensuremath{\checkmark}

\textbf{تفسیر.} هر سه رادیولوژیست در عمل حساسیت یکسانی دارند
(\(d' \approx 2\)). \lr{R1} و \lr{R2} درصد درست \emph{یکسان} ولی سوگیری
\emph{مخالف} دارند. اینکه کدام سوگیری بهتر است سؤالی درباره‌ی ادراک نیست،
بلکه درباره‌ی هزینه‌هاست: از دست دادن یک تومور معمولاً خیلی بدتر از یک
\lr{false alarm} است، پس شاید ملاک \lr{liberal} همان ملاک \emph{مطلوب}
باشد. \lr{SDT} به ما اجازه می‌دهد «آیا می‌بینند؟» و «چقدر آماده‌اند
بگویند؟» را جداگانه بپرسیم.

\textbf{تصحیح نرخ حدی.} یک شرکت‌کننده 50/50 \lr{hit} و 3/50 \lr{false
alarm} دارد. پس \(H = 1\) و \(d' = \infty\). با تصحیح \lr{log-linear:}
\(H = 50.5/51 = 0.990\)، \(F = 3.5/51 = 0.069\)، پس
\(d' = 2.33 - (-1.49) = 3.82\).

\subsubsection{بخش 2: فراموشی نمایی در برابر توانی}

برای هر دو مدل \(A = 0.95\) بگیرید، برای مدل نمایی \(\lambda = 0.25\) در
روز و برای مدل توانی \(d = 0.6\).

{\def\LTcaptype{none} % do not increment counter
\begin{landscape}\small
\begin{longtable}[]{@{}
  >{\raggedright\arraybackslash}p{(\linewidth - 12\tabcolsep) * \real{0.1429}}
  >{\raggedright\arraybackslash}p{(\linewidth - 12\tabcolsep) * \real{0.1429}}
  >{\raggedright\arraybackslash}p{(\linewidth - 12\tabcolsep) * \real{0.1429}}
  >{\raggedright\arraybackslash}p{(\linewidth - 12\tabcolsep) * \real{0.1429}}
  >{\raggedright\arraybackslash}p{(\linewidth - 12\tabcolsep) * \real{0.1429}}
  >{\raggedright\arraybackslash}p{(\linewidth - 12\tabcolsep) * \real{0.1429}}
  >{\raggedright\arraybackslash}p{(\linewidth - 12\tabcolsep) * \real{0.1429}}@{}}
\toprule\noalign{}
\begin{minipage}[b]{\linewidth}\raggedright
\(t\) (روز)
\end{minipage} & \begin{minipage}[b]{\linewidth}\raggedright
0
\end{minipage} & \begin{minipage}[b]{\linewidth}\raggedright
1
\end{minipage} & \begin{minipage}[b]{\linewidth}\raggedright
2
\end{minipage} & \begin{minipage}[b]{\linewidth}\raggedright
7
\end{minipage} & \begin{minipage}[b]{\linewidth}\raggedright
14
\end{minipage} & \begin{minipage}[b]{\linewidth}\raggedright
30
\end{minipage} \\
\midrule\noalign{}
\endhead
\bottomrule\noalign{}
\endlastfoot
نمایی \(0.95\,e^{-0.25t}\) & 0.950 & 0.740 & 0.576 & 0.165 & 0.029 &
0.001 \\
توانی \(0.95\,(1+t)^{-0.6}\) & 0.950 & 0.627 & 0.491 & 0.273 & 0.187 &
0.121 \\
\end{longtable}
\end{landscape}
}

\textbf{نیمه‌عمر (زمانی که یادآوری به نصف سطح اولیه می‌رسد):} - نمایی:
\(e^{-\lambda t} = 1/2 \Rightarrow t_{1/2} = \ln 2 / \lambda = 0.693/0.25 = 2.77\)
روز. - توانی:
\((1+t)^{-d} = 1/2 \Rightarrow t_{1/2} = 2^{1/d} - 1 = 2^{1.667} - 1 = 2.17\)
روز.

\textbf{تفسیر.} در دو روز اول منحنی‌ها نسبتاً شبیه‌اند: حدود 0.1 با هم فرق
دارند که با \(n = 50\) آیتم فقط حدود 1.5 انحراف معیار دوجمله‌ای است. بعد
از یک ماه کاملاً با هم اختلاف دارند: مدل نمایی می‌گوید تقریباً همه چیز
فراموش شده (0.1\%)، در حالی که مدل توانی می‌گوید 12\% هنوز باقی است.
\textbf{فاصله‌های نگهداری طولانی جایی است که این مدل‌ها از هم تشخیص داده
می‌شوند}؛ درسی درباره‌ی طراحی آزمایش که در تکلیف \lr{A1} از آن استفاده
می‌کنید.

\begin{center}\rule{0.5\linewidth}{0.5pt}\end{center}

\subsection{\texorpdfstring{\lr{lab} (2 ساعت)}{lab (2 ساعت)}}

\subsubsection{اهداف}

\begin{enumerate}
\def\labelenumi{\arabic{enumi}.}
\tightlist
\item
  پیاده‌سازی \lr{SDT} به‌عنوان یک مدل مولد \lr{(generative)}، شبیه‌سازی
  آزمایش‌های بله/نه، و برآورد \(d'\) و \(c\).
\item
  استفاده از نمودارهای \lr{ROC} و \lr{zROC} برای \textbf{ارزیابی} فرض
  واریانس برابر.
\item
  پیاده‌سازی فراموشی نمایی و توانی با نویز دوجمله‌ای و بررسی
  پارامترهایشان.
\item
  \textbf{برازش} هر دو مدل فراموشی با خطی‌سازی، \textbf{ارزیابی} آن‌ها با
  باقی‌مانده‌ها \lr{(residuals)}، و \textbf{تفسیر} نتیجه.
\end{enumerate}

\textbf{زمان‌بندی:} بخش \lr{A} 30 دقیقه · بخش \lr{B} 25 دقیقه · بخش
\lr{C} 30 دقیقه · بخش \lr{D} 25 دقیقه · زمان باقی‌مانده: تمرین‌ها.

\subsubsection{آماده‌سازی}

\begin{quote}
\textbf{\lr{Code} · کد:} \lr{\nolinkurl{code/week02.py}}, \lr{cell} 1 ·
سلول 1
\end{quote}

\subsubsection{\texorpdfstring{بخش \lr{A: SDT} به‌عنوان یک مدل مولد (30
دقیقه)}{بخش A: SDT به‌عنوان یک مدل مولد (30 دقیقه)}}

\begin{quote}
\textbf{\lr{Code} · کد:} \lr{\nolinkurl{code/week02.py}}, \lr{cell} 2 ·
سلول 2
\end{quote}

یک شرکت‌کننده را شبیه‌سازی کنید و پارامترها را بازیابی کنید:

\begin{quote}
\textbf{\lr{Code} · کد:} \lr{\nolinkurl{code/week02.py}}, \lr{cell} 3 ·
سلول 3
\end{quote}

\textbf{خروجی مورد انتظار:}

\Needspace{8\baselineskip}
\begin{latin}
\begin{Verbatim}
Counts: {'hits': 68, 'misses': 32, 'fas': 11, 'crs': 89}
Predicted H = 0.674, F = 0.147
Estimated d' = 1.67 (true 1.5),  c = 0.37 (true 0.3)
\end{Verbatim}
\end{latin}

یک شرکت‌کننده‌ی شبیه‌سازی‌شده چیز زیادی درباره‌ی کیفیت \emph{برآوردگر}
نمی‌گوید. آزمایش را چندین بار تکرار کنیم: این پیش‌نمایشی از منطق
\textbf{بازیابی پارامتر} در هفته‌ی 7 است.

\begin{quote}
\textbf{\lr{Code} · کد:} \lr{\nolinkurl{code/week02.py}}, \lr{cell} 4 ·
سلول 4
\end{quote}

\textbf{خروجی مورد انتظار:}

\Needspace{11\baselineskip}
\begin{latin}
\begin{Verbatim}
            d_hat         c_hat
             mean    std   mean    std
n_per_type
20          1.500  0.449  0.296  0.224
50          1.489  0.290  0.295  0.143
200         1.502  0.142  0.299  0.071
\end{Verbatim}
\end{latin}

\textbf{به این دقت کنید:} برآوردها حول مقادیر واقعی پراکنده‌اند و
پراکندگی با افزایش تعداد \lr{trial}ها کم می‌شود. با 20 \lr{trial} از هر
نوع، انحراف معیار \(\hat d'\) حدود 0.45 است، پس برآوردهای تکی اغلب \ensuremath{\pm}0.5
یا بیشتر خطا دارند، و برای نصف کردن پراکندگی تقریباً 4 برابر \lr{trial}
لازم است. اینجا به‌طور میانگین برآوردها درست سر هدف‌اند. ولی با \(d'\)
بالاتر و \lr{trial}های کم، نرخ‌های حدی رایج می‌شوند و آن‌وقت تصحیح،
برآوردها را به سمت 0 می‌کشد (\lr{\nolinkurl{true_d = 3}} را با 20
\lr{trial} امتحان کنید). هر وقت \(d'\) گروه‌هایی را مقایسه می‌کنید که با
تعداد \lr{trial} متفاوت آزمون شده‌اند، این را در نظر داشته باشید.

\subsubsection{\texorpdfstring{بخش \lr{B:} منحنی‌های \lr{ROC} و فرض
واریانس برابر (25
دقیقه)}{بخش B: منحنی‌های ROC و فرض واریانس برابر (25 دقیقه)}}

اول، \lr{ROC}های پیش‌بینی‌شده‌ی مدل:

\begin{quote}
\textbf{\lr{Code} · کد:} \lr{\nolinkurl{code/week02.py}}, \lr{cell} 5 ·
سلول 5
\end{quote}

حالا یک آزمایش \textbf{درجه‌بندی اطمینان} را شبیه‌سازی کنید. پاسخ‌ها روی یک
مقیاس 6 درجه‌ای داده می‌شوند (1 = «مطمئنم نویز است» \ensuremath{\ldots} 6 = «مطمئنم
علامت است») که معادل پنج ملاک هم‌زمان است. داده‌ها را از یک مشاهده‌گر با
واریانس برابر و از یک مشاهده‌گر با \textbf{واریانس نابرابر} که توزیع
علامتش انحراف معیار 1.5 دارد تولید می‌کنیم.

\begin{quote}
\textbf{\lr{Code} · کد:} \lr{\nolinkurl{code/week02.py}}, \lr{cell} 6 ·
سلول 6
\end{quote}

\textbf{خروجی مورد انتظار:}

\Needspace{7\baselineskip}
\begin{latin}
\begin{Verbatim}
sigma_signal = 1.0: zROC slope = 0.96 (predicted 1.00); d' at each criterion = [1.52 1.48 1.48 1.43 1.43]
sigma_signal = 1.5: zROC slope = 0.67 (predicted 0.67); d' at each criterion = [1.46 1.34 1.16 0.96 0.84]
\end{Verbatim}
\end{latin}

\textbf{ارزیابی و تفسیر.} برای مشاهده‌گر با واریانس برابر، شیب \lr{zROC}
نزدیک 1 است و \(d'\) در همه‌ی ملاک‌ها تقریباً یکسان است: فرض مدل برقرار
است. برای مشاهده‌گر با واریانس نابرابر، شیب نزدیک \(1/1.5 = 0.67\) است و
\(d'\) \textbf{به‌طور نظام‌مند با ملاک تغییر می‌کند}. اگر از چنین مشاهده‌گری
فقط داده‌ی بله/نه جمع کرده بودیم، تفاوت گروه‌ها در «\(d'\)» می‌توانست صرفاً
از تفاوت در سوگیری بیاید. بررسی پیش‌بینی امضایی یک مدل (اینجا: شیب
\lr{zROC} = 1) شکل پایه‌ای از \textbf{ارزیابی مدل} است.

\subsubsection{\texorpdfstring{بخش \lr{C:} منحنی‌های فراموشی با نویز
دوجمله‌ای (30
دقیقه)}{بخش C: منحنی‌های فراموشی با نویز دوجمله‌ای (30 دقیقه)}}

\begin{quote}
\textbf{\lr{Code} · کد:} \lr{\nolinkurl{code/week02.py}}, \lr{cell} 7 ·
سلول 7
\end{quote}

پارامترها را بررسی کنید: \(A\) و پارامتر زوال هر منحنی را چگونه تغییر
می‌دهند؟

\begin{quote}
\textbf{\lr{Code} · کد:} \lr{\nolinkurl{code/week02.py}}, \lr{cell} 8 ·
سلول 8
\end{quote}

حالا یک \textbf{آزمایش شبیه‌سازی‌شده} بسازید. داده‌ها از مدل \emph{توانی}
می‌آیند، با 50 آیتم که در هر یک از 8 فاصله‌ی نگهداری آزمون می‌شوند:

\begin{quote}
\textbf{\lr{Code} · کد:} \lr{\nolinkurl{code/week02.py}}, \lr{cell} 9 ·
سلول 9
\end{quote}

\textbf{خروجی مورد انتظار:}

\Needspace{14\baselineskip}
\begin{latin}
\begin{Verbatim}
       t  recalled   n  p_obs
0   0.25        42  50   0.84
1   1.00        30  50   0.60
2   2.00        26  50   0.52
3   4.00        25  50   0.50
4   7.00        13  50   0.26
5  14.00        13  50   0.26
6  21.00        13  50   0.26
7  30.00        14  50   0.28
\end{Verbatim}
\end{latin}

\subsubsection{\texorpdfstring{بخش \lr{D:} برازش با خطی‌سازی، ارزیابی،
تفسیر (25 دقیقه)}{بخش D: برازش با خطی‌سازی، ارزیابی، تفسیر (25 دقیقه)}}

\textbf{برازش.} هر مدل را با رگرسیون خط راست روی مقیاس لگاریتمی برازش
کنید (بخش 5 جزوه).

\begin{quote}
\textbf{\lr{Code} · کد:} \lr{\nolinkurl{code/week02.py}}, \lr{cell} 10 ·
سلول 10
\end{quote}

\textbf{ارزیابی.} به منحنی‌های برازش‌شده \emph{و} باقی‌مانده‌ها نگاه کنید.
یک مدل خوب باقی‌مانده‌هایی به جا می‌گذارد که به‌طور تصادفی حول صفر
پراکنده‌اند؛ یک الگوی نظام‌مند یعنی \emph{شکل} مدل غلط است.

\begin{quote}
\textbf{\lr{Code} · کد:} \lr{\nolinkurl{code/week02.py}}, \lr{cell} 11 ·
سلول 11
\end{quote}

\textbf{خروجی مورد انتظار:}

\Needspace{8\baselineskip}
\begin{latin}
\begin{Verbatim}
exp: A = 0.550, rate = 0.032, RMSD = 0.1308
pow: A = 0.797, rate = 0.370, RMSD = 0.0628
true (power): A = 0.9, d = 0.5
\end{Verbatim}
\end{latin}

\textbf{تفسیر (برای هرکدام 2 تا 3 جمله در \lr{notebook} بنویسید):}

\begin{enumerate}
\def\labelenumi{\arabic{enumi}.}
\tightlist
\item
  کدام مدل داده‌ها را بهتر توصیف می‌کند، و الگوی باقی‌مانده‌های مدل
  \emph{دیگر} چه شکلی دارد؟ (دنبال شکل \lr{U} یا \lr{U} وارونه بگردید:
  مدل نمایی در فاصله‌های میانی بیش‌برآورد و در فاصله‌های طولانی کم‌برآورد
  می‌کند، یا برعکس.)
\item
  با پارامترهای برازش‌شده، نرخ نسبی فراموشی هر مدل بعد از 1 روز و بعد از
  20 روز چقدر است؟ این برای حافظه چه معنایی دارد؟
\item
  اینجا مدل توانی مدل واقعی است، ولی برآوردهایش
  (\(\hat A \approx 0.80\)، \(\hat d \approx 0.37\)) آشکارا با مقادیر
  واقعی (\(A^\ast = 0.9\)، \(d^\ast = 0.5\)) فاصله دارند. دو دلیل نام
  ببرید. (راهنمایی: به سه نقطه‌ی آخر نگاه کنید؛ فکر کنید تبدیل لگاریتمی
  با وزن نسبت‌های کوچک چه می‌کند.) در هفته‌ی 3 همین داده‌ها را روی مقیاس
  اصلی دوباره برازش و مقایسه می‌کنیم.
\end{enumerate}

\subsubsection{تمرین‌ها (به ترتیب سختی)}

\textbf{تمرین 1 \ensuremath{\star} (نرخ‌های حدی).} یک شرکت‌کننده 40/40 \lr{hit} و 0/40
\lr{false alarm} دارد. \(d'\) را بدون تصحیح و با تصحیح \lr{log-linear}
حساب کنید. \(d'\) تصحیح‌شده نماینده‌ی چیست، و چرا باید تصحیحی را که
استفاده کرده‌اید گزارش کنید؟

\textbf{تمرین 2 \ensuremath{\star} (نیمه‌عمرها).} با پارامترهای برازش‌شده از بخش \lr{D}،
نیمه‌عمر هر مدل را حساب کنید. بعد حساب کنید چقدر طول می‌کشد تا یادآوری به
10\% سطح اولیه برسد. مدل‌ها کجا بیشترین اختلاف را دارند؟

\textbf{تمرین 3 \ensuremath{\star}\ensuremath{\star} (داده‌های ابینگهاوس).} داده‌های کلاسیک \lr{savings} از
آزمایش‌های ابینگهاوس روی خودش در 1885 اغلب این‌طور گزارش می‌شوند {[}پیش از
استفاده در کلاس با منبع اصلی تطبیق داده شود{]}:

{\def\LTcaptype{none} % do not increment counter
\begin{landscape}\small
\begin{longtable}[]{@{}
  >{\raggedright\arraybackslash}p{(\linewidth - 14\tabcolsep) * \real{0.1250}}
  >{\raggedright\arraybackslash}p{(\linewidth - 14\tabcolsep) * \real{0.1250}}
  >{\raggedright\arraybackslash}p{(\linewidth - 14\tabcolsep) * \real{0.1250}}
  >{\raggedright\arraybackslash}p{(\linewidth - 14\tabcolsep) * \real{0.1250}}
  >{\raggedright\arraybackslash}p{(\linewidth - 14\tabcolsep) * \real{0.1250}}
  >{\raggedright\arraybackslash}p{(\linewidth - 14\tabcolsep) * \real{0.1250}}
  >{\raggedright\arraybackslash}p{(\linewidth - 14\tabcolsep) * \real{0.1250}}
  >{\raggedright\arraybackslash}p{(\linewidth - 14\tabcolsep) * \real{0.1250}}@{}}
\toprule\noalign{}
\begin{minipage}[b]{\linewidth}\raggedright
فاصله
\end{minipage} & \begin{minipage}[b]{\linewidth}\raggedright
20 دقیقه
\end{minipage} & \begin{minipage}[b]{\linewidth}\raggedright
1 ساعت
\end{minipage} & \begin{minipage}[b]{\linewidth}\raggedright
8.8 ساعت
\end{minipage} & \begin{minipage}[b]{\linewidth}\raggedright
1 روز
\end{minipage} & \begin{minipage}[b]{\linewidth}\raggedright
2 روز
\end{minipage} & \begin{minipage}[b]{\linewidth}\raggedright
6 روز
\end{minipage} & \begin{minipage}[b]{\linewidth}\raggedright
31 روز
\end{minipage} \\
\midrule\noalign{}
\endhead
\bottomrule\noalign{}
\endlastfoot
\lr{Savings} (\%) & 58.2 & 44.2 & 35.8 & 33.7 & 27.8 & 25.4 & 21.1 \\
\end{longtable}
\end{landscape}
}

فاصله‌ها را به روز تبدیل کنید، هر دو مدل را با خطی‌سازی برازش کنید و
باقی‌مانده‌هایشان را مقایسه کنید. کدام شکل بهتر برازش می‌شود؟ بعد تحلیل را
با زمان بر حسب \textbf{ساعت} تکرار کنید. آیا برازش توانی تغییر می‌کند؟ با
استفاده از کادر «پیشرفته» در بخش 3 جزوه توضیح دهید.

\textbf{تمرین 4 \ensuremath{\star}\ensuremath{\star} (کدام طرح آزمایش مدل‌ها را از هم جدا می‌کند؟).}
\((A, \lambda)\) و \((A, d)\) را طوری انتخاب کنید که دو مدل در \(t = 1\)
و \(t = 7\) روز یادآوری \emph{یکسانی} پیش‌بینی کنند. (راهنمایی: برای هر
مدل دو معادله و دو مجهول؛ لگاریتم بگیرید.) بعد هر دو منحنی را از
\(t = 0\) تا \(60\) روز رسم کنید و جاهایی را که اختلافشان با \(n = 50\)
از 2 انحراف معیار دوجمله‌ای بیشتر است علامت بزنید. کدام فاصله‌های نگهداری
را آزمون می‌کردید؟

\textbf{تمرین 5 \ensuremath{\star}\ensuremath{\star}\ensuremath{\star} (پیشرفته: شمارش صفر).} از مدل نمایی با \(A = 0.9\)،
\(\lambda = 0.3\) در فاصله‌های بخش \lr{C} و با \(n = 20\) آیتم داده تولید
کنید. چند بار (از 1000 مجموعه‌داده‌ی شبیه‌سازی‌شده) دست‌کم یک فاصله یادآوری
صفر دارد که \lr{\nolinkurl{fit_linearized}} را از کار می‌اندازد؟ یک راه‌حل
سریع (مثلاً افزودن 0.5 به شمارش‌ها) پیشنهاد و پیاده‌سازی کنید و بررسی کنید
آیا برآورد \(\lambda\) را سوگیر می‌کند. توضیح دهید چرا این نشان می‌دهد به
روشی نیاز داریم که روی مقیاس اصلی کار کند (هفته‌های 3\lr{--}4).

\subsubsection{\texorpdfstring{\lr{lab check-off}}{lab check-off}}

به مدرس یا دستیار آموزشی نشان دهید: (1) جدول بازیابی \lr{SDT} از بخش
\lr{A}، (2) دو \lr{zROC} تجربی با شیب‌هایشان از بخش \lr{B}، و (3) نمودار
باقی‌مانده‌ها از بخش \lr{D} با یک جمله تفسیر.

\begin{center}\rule{0.5\linewidth}{0.5pt}\end{center}

\subsection{\texorpdfstring{ارتباط با هوش مصنوعی و
\lr{LLM}ها}{ارتباط با هوش مصنوعی و LLMها}}

\textbf{آشکارسازهای متن هوش مصنوعی، مشاهده‌گرهایی در یک تکلیف \lr{SDT}
هستند.} دانشگاه‌ها روزبه‌روز بیشتر از ابزارهایی استفاده می‌کنند که به یک
متن امتیاز «چقدر محتمل است که هوش مصنوعی نوشته باشد» می‌دهند. هر متن یا
با هوش مصنوعی نوشته شده («علامت») یا با انسان («نویز»)، و ابزار وقتی
امتیاز از یک آستانه بیشتر شود می‌گوید «هوش مصنوعی». این دقیقاً موقعیت
\lr{SDT} این هفته است: ابزار یک \textbf{حساسیت} \(d'\) دارد (امتیاز چقدر
خوب دو نوع متن را از هم جدا می‌کند) و یک \textbf{ملاک} \(c\) (آستانه‌ای که
کاربر انتخاب می‌کند). یک \lr{false alarm} یعنی دانشجویی به‌اشتباه متهم
شود، پس انتخاب ملاک یک تصمیم اخلاقی است، نه فقط فنی.

\begin{quote}
\textbf{\lr{Code} · کد:} \lr{\nolinkurl{code/week02.py}}, \lr{cell} 12 ·
سلول 12
\end{quote}

\textbf{خروجی مورد انتظار:}

\Needspace{8\baselineskip}
\begin{latin}
\begin{Verbatim}
threshold 0.0: H = 0.84, F = 0.49, d' = 1.03, c = -0.50
threshold 0.6: H = 0.70, F = 0.28, d' = 1.11, c = 0.03
threshold 1.2: H = 0.51, F = 0.09, d' = 1.35, c = 0.66
\end{Verbatim}
\end{latin}

جابه‌جا کردن آستانه \(c\) را خیلی تغییر می‌دهد (از \ensuremath{-}0.50 تا 0.66)، در حالی
که \(d'\) نزدیک مقدار واقعی‌اش یعنی 1.2 می‌ماند. تغییرات \(d'\) نویز
نمونه‌گیری ناشی از 200 متن در هر گروه است.

\textbf{امتحانش کنید (بدون نیاز به \lr{API).}} - \ensuremath{\star} فرض کنید 10\%
مقاله‌های تحویلی با هوش مصنوعی نوشته شده‌اند. در آستانه‌ی 1.2،
\(P(\text{\lr{AI-written}} \mid \text{\lr{flagged}})\) چقدر است؟ (از قاعده‌ی بیز با
\(H\) و \(F\) استفاده کنید.) آیا بر اساس یک هشدار اقدام می‌کنید؟ - \ensuremath{\star}\ensuremath{\star}
\lr{ROC} آشکارساز (همه‌ی آستانه‌ها از \ensuremath{-}2 تا 3) و \lr{zROC} آن را رسم کنید.
آیا شیب نزدیک 1 است؟ با امتیازهای متن انسانی با انحراف معیار 1.5
(آشکارسازی با واریانس نابرابر) تکرار کنید. - \ensuremath{\star}\ensuremath{\star} بحث: گزارش شده که
آشکارسازها متن‌های نویسندگان غیربومی انگلیسی را بیشتر علامت می‌زنند
{[}بررسی شود: \lr{Liang et al.,} 2023, \emph{\lr{Patterns}}{]}. به زبان
\lr{SDT} این یعنی چه، و یک دانشگاه چه باید بکند؟

\begin{center}\rule{0.5\linewidth}{0.5pt}\end{center}

\subsection{اشتباهات رایج و برداشت‌های نادرست}

\begin{enumerate}
\def\labelenumi{\arabic{enumi}.}
\tightlist
\item
  \textbf{«درصد پاسخ درست کافی است.»} حساسیت و سوگیری را قاطی می‌کند
  (رادیولوژیست‌ها را ببینید). هر وقت پاسخ‌ها ممکن است سوگیری داشته باشند
  از \(d'\) و \(c\) استفاده کنید.
\item
  \textbf{محاسبه‌ی \(d'\) به شکل \(\Phi^{-1}(H - F)\) یا از روی درصدها.}
  فرمول \(\Phi^{-1}(H) - \Phi^{-1}(F)\) است، با نرخ‌ها به‌صورت نسبت بین 0
  و 1.
\item
  \textbf{فراموش کردن مشکل نرخ‌های حدی.} \(H = 1\) یا \(F = 0\) یک \(d'\)
  بی‌نهایت می‌دهد و \lr{Python} بدون هیچ خطایی \lr{\nolinkurl{inf}}
  برمی‌گرداند. همیشه از یک تصحیح استفاده کنید (و آن را گزارش کنید).
\item
  \textbf{فرض اینکه \(d'\) در هر حالتی مستقل از سوگیری است.} فقط تحت مدل
  واریانس برابر مستقل از ملاک است. اگر شیب \lr{zROC} با 1 فرق کند،
  \(d'\) یک جفت \((H, F)\) به سوگیری بستگی دارد.
\item
  \textbf{«نمایی و توانی فقط دو شکل منحنی‌اند.»} آن‌ها ادعاهای متفاوتی
  درباره‌ی چگونگی فراموشی دارند (نرخ نسبی ثابت در برابر کاهشی).
\item
  \textbf{قضاوت درباره‌ی مدل فقط با آماره‌ی برازش.} \lr{RMSD} کوچک همراه
  با الگوی نظام‌مند در باقی‌مانده‌ها یعنی شکل مدل غلط است. همیشه
  باقی‌مانده‌ها را رسم کنید.
\item
  \textbf{نادیده گرفتن مدل نویز.} مدل فقط \(f(x;\boldsymbol\theta)\)
  نیست. نویز دوجمله‌ای نسبت‌های یادآوری جزئی از مدل است و به شما می‌گوید
  کدام تفاوت‌ها معنادارند.
\item
  \textbf{تغییر واحد \(t\) در مدل توانی جابه‌جاشده} و انتظار شکل یکسان.
\end{enumerate}

\begin{center}\rule{0.5\linewidth}{0.5pt}\end{center}

\subsection{خودآزمایی}

\textbf{1.} پنج جزء یک مدل را نام ببرید و نسخه‌ی \lr{SDT} هرکدام را
بگویید. \textbf{پاسخ} طرح آزمایش \lr{(trial}های علامت در برابر نویز؛
پاسخ‌های بله/نه)؛ پارامترها (آزاد: \(d'\)، \(c\)؛ ثابت: میانگین نویز 0،
هر دو انحراف معیار 1)؛ هسته‌ی قطعی (\(H = \Phi(d'/2 - c)\)،
\(F = \Phi(-d'/2 - c)\))؛ مدل نویز (شمارش \lr{hit}ها و \lr{false
alarm}ها دوجمله‌ای است)؛ فرض پیوند (وقتی شواهد از ملاک بیشتر است پاسخ
«بله»).

\textbf{2.} چرا میانگین نویز و انحراف معیارها در \lr{SDT} ثابت‌اند و آزاد
نیستند؟ \textbf{پاسخ} چون \lr{identifiable} نیستند: جابه‌جا کردن هر دو
توزیع و ملاک به یک اندازه، یا کشیدن کل محور، هیچ پیش‌بینی‌ای را عوض
نمی‌کند. ثابت کردن آن‌ها (قیدهای مقیاس) واحدهایی را تعریف می‌کند که \(d'\)
و \(c\) با آن اندازه‌گیری می‌شوند.

\textbf{3.} یک مشاهده‌گر \(H = 0.80\) و \(F = 0.20\) دارد. \(d'\) و \(c\)
را حساب کنید. (\(\Phi^{-1}(0.80) = 0.842\)) \textbf{پاسخ}
\(d' = 0.842 - (-0.842) = 1.68\)؛
\(c = -\tfrac{1}{2}(0.842 - 0.842) = 0\)، پس مشاهده‌گر بی‌سوگیری است.

\textbf{4.} \lr{SDT} با واریانس برابر درباره‌ی \lr{zROC} چه پیش‌بینی
می‌کند، و اگر شیب تجربی 0.7 باشد یعنی چه؟ \textbf{پاسخ} یک خط راست با شیب
1. شیب 0.7 نشان می‌دهد توزیع علامت از توزیع نویز متغیرتر است
(\(\sigma_s \approx 1/0.7 \approx 1.4\)). در این صورت \(d'\) تک‌نقطه‌ای به
ملاک بستگی دارد و باید از مدل واریانس نابرابر (یا \(d_a\)) استفاده کرد.

\textbf{5.} نرخ نسبی فراموشی هر مدل چیست، و چرا از نظر روان‌شناختی مهم
است؟ \textbf{پاسخ} نمایی: \(\lambda\)، ثابت. توانی: \(d/(1+t)\)، با گذشت
زمان کاهشی. مدل نمایی می‌گوید خاطره‌ها در هر سنی روزانه نسبت یکسانی از
قدرتشان را از دست می‌دهند؛ تابع توانی می‌گوید خاطره‌های قدیمی‌تر در برابر
فراموشی بیشتر مقاوم‌تر می‌شوند.

\textbf{6.} برای اینکه سریع ببینید داده‌های فراموشی نمایی به نظر می‌رسند
چه نموداری رسم می‌کنید؟ و برای توانی؟ \textbf{پاسخ} نمایی:
\(\ln(\text{\lr{recall}})\) در برابر \(t\) (نیمه‌لگاریتمی) که باید خط راست
باشد. توانی: \(\ln(\text{\lr{recall}})\) در برابر \(\ln(1+t)\)
(لگاریتمی\lr{--}لگاریتمی) که باید خط راست باشد.

\textbf{7.} دو دلیل بیاورید که چرا برازش خط راست روی نسبت‌های لگاریتمی‌شده
روش برآورد ضعیفی است. \textbf{پاسخ} (1) وقتی یک نسبت مشاهده‌شده 0 باشد از
کار می‌افتد. (2) نویز را تحریف می‌کند: نسبت‌های کوچک روی مقیاس لگاریتمی
خیلی متغیر می‌شوند ولی با وزن برابر در نظر گرفته می‌شوند، و خطاهایی که
کمینه می‌شوند روی مقیاس غلط (لگاریتمی) هستند. برازش روی مقیاس اصلی با
کمترین مربعات یا بیشینه درست‌نمایی از این مشکلات جلوگیری می‌کند.

\begin{center}\rule{0.5\linewidth}{0.5pt}\end{center}

\subsection{تکلیف}

\subsubsection{\texorpdfstring{\lr{A1:} از کلامی به فرمال (1.0 نمره از
20؛ انتشار هفته‌ی 2، موعد ابتدای درس هفته‌ی
4)}{A1: از کلامی به فرمال (1.0 نمره از 20؛ انتشار هفته‌ی 2، موعد ابتدای درس هفته‌ی 4)}}

\textbf{سناریو.} دو گروه پژوهشی فراموشی واژه‌هایی را که در یک کلاس زبان
خارجی یاد گرفته شده‌اند مطالعه می‌کنند. - \textbf{گروه \lr{C}} می‌گوید:
\emph{«خاطره‌ها با نرخ ثابتی زوال پیدا می‌کنند.»} - \textbf{گروه \lr{R}}
می‌گوید: \emph{«خاطره‌ها با بالا رفتن سن در برابر فراموشی مقاوم‌تر
می‌شوند.»}

هر دو گروه ادعا می‌کنند داده‌های منتشرشده (یادآوری سریع افت می‌کند و بعد
ثابت می‌شود) از دیدگاهشان حمایت می‌کند.

\textbf{کارها.} یک \lr{Jupyter notebook} به‌همراه خروجی \lr{PDF} آن تحویل
دهید.

\begin{enumerate}
\def\labelenumi{\arabic{enumi}.}
\tightlist
\item
  \textbf{فرمال‌سازی (گروه \lr{C} و گروه \lr{R).}} هر ادعای کلامی را
  به‌صورت یک مدل فرمال با هر پنج جزء بخش 1 جزوه بنویسید: طرح آزمایش،
  پارامترهای آزاد/ثابت با دامنه‌شان، هسته‌ی قطعی، مدل نویز دوجمله‌ای، و فرض
  پیوند. در 2 تا 3 جمله توجیه کنید چرا هر معادله ادعای کلامی‌اش را نشان
  می‌دهد (از نرخ نسبی فراموشی استفاده کنید).
\item
  \textbf{پیاده‌سازی.} توابع \lr{Python} بنویسید که (الف) احتمال‌های
  یادآوری پیش‌بینی‌شده را برگردانند و (ب) برای فاصله‌ها و تعداد آیتم \(n\)
  داده‌شده، شمارش‌های یادآوری دوجمله‌ای را شبیه‌سازی کنند.
\item
  \textbf{کالیبراسیون.} هر دو گروه قبول دارند که یادآوری \textbf{بعد از
  1 روز 0.60} و \textbf{بعد از 7 روز 0.35} است. پارامترهای هر مدل را که
  دقیقاً این دو مقدار را بازتولید می‌کنند، دستی و به‌صورت بسته به دست
  آورید. محاسبه را نشان دهید (لگاریتم بگیرید) و آن را عددی بررسی کنید.
\item
  \textbf{پیدا کردن جای اختلاف.} با پارامترهای کالیبره‌شده، هر دو منحنی
  پیش‌بینی‌شده را از 0 تا 90 روز رسم کنید. برای \(n = 40\) آیتم، در هر
  فاصله یک \emph{شاخص تمایزپذیری}
  \(D(t) = \big(f_C(t) - f_R(t)\big)^2 / \mathrm{Var}(t)\) حساب کنید که
  در آن \(\mathrm{Var}(t) = \bar f(1 - \bar f)/n\) و \(\bar f\) میانگین
  دو پیش‌بینی است. \(D(t)\) را رسم کنید.
\item
  \textbf{طراحی آزمایش.} می‌توانید هر شرکت‌کننده را فقط در \textbf{4}
  فاصله‌ی نگهداری آزمون کنید که از میان \{0.5، 1، 2، 4، 7، 14، 30، 60،
  90\} روز انتخاب می‌شوند. مجموعه‌ای را انتخاب کنید که مجموع \(D(t)\) را
  بیشینه کند، توضیح دهید چرا منطقی است، و با طرح خود یک مجموعه‌داده از
  \emph{هر} مدل شبیه‌سازی کنید. دو مجموعه‌داده را روی یک نمودار نشان دهید.
\item
  \textbf{تفسیر (حداکثر 300 کلمه).} کدام داده‌ها به نفع گروه \lr{C} و
  کدام به نفع گروه \lr{R} هستند؟ تحلیل شما درباره‌ی اینکه چرا هر دو گروه
  می‌توانند ادعا کنند «داده‌های منتشرشده» از آن‌ها حمایت می‌کند چه می‌گوید؟
  یک فرض مدل‌هایتان را که ممکن است غلط باشد نام ببرید و بگویید چطور
  نتیجه‌گیری‌تان را تغییر می‌دهد.
\end{enumerate}

\textbf{روبریک نمره‌دهی (10 امتیاز، معادل 1.0 از نمره‌ی نهایی):}

{\def\LTcaptype{none} % do not increment counter
\begin{longtable}[]{@{}
  >{\raggedright\arraybackslash}p{(\linewidth - 4\tabcolsep) * \real{0.3333}}
  >{\raggedright\arraybackslash}p{(\linewidth - 4\tabcolsep) * \real{0.3333}}
  >{\raggedright\arraybackslash}p{(\linewidth - 4\tabcolsep) * \real{0.3333}}@{}}
\toprule\noalign{}
\begin{minipage}[b]{\linewidth}\raggedright
معیار
\end{minipage} & \begin{minipage}[b]{\linewidth}\raggedright
امتیاز
\end{minipage} & \begin{minipage}[b]{\linewidth}\raggedright
شرط نمره‌ی کامل
\end{minipage} \\
\midrule\noalign{}
\endhead
\bottomrule\noalign{}
\endlastfoot
فرمال‌سازی (کار 1) & 2 & هر پنج جزء برای هر دو مدل؛ دامنه‌ی درست
پارامترها؛ ارتباط معادله و ادعای کلامی با نرخ نسبی فراموشی توضیح داده
شده \\
کد (کار 2) & 2 & درست، با توضیح \lr{(comment)}، قابل بازتولید \lr{(seed}
ثابت)؛ توابع برداری‌شده یا با حلقه‌ی روشن \\
کالیبراسیون (کار 3) & 2 & استخراج بسته‌ی درست و بررسی عددی برای هر دو
مدل \\
اختلاف و طراحی (کارهای 4\lr{--}5) & 2.5 & \(D(t)\) درست؛ طرح انتخاب‌شده توجیه
شده و واقعاً معیار را بیشینه می‌کند؛ نمودارهای روشن \\
تفسیر (کار 6) & 1.5 & روشن، مشخص، با ارجاع به نتایج خودتان؛ یک محدودیت
واقعی را نام می‌برد \\
\end{longtable}
}

برای تحویل با تأخیر، روزی 10\% کسر می‌شود (حداکثر 3 روز).

\begin{center}\rule{0.5\linewidth}{0.5pt}\end{center}

\subsection{مطالعه‌ی بیشتر}

\begin{itemize}
\tightlist
\item
  \textbf{\lr{Farrell} \lr{\&} \lr{Lewandowsky} (2018)، فصل 2.} ساختن مدل از
  یک نظریه‌ی کلامی؛ اجزای مدل و نقش پارامترها. \emph{(منبع اصلی.)}
\item
  \textbf{\lr{Lee} \lr{\&} \lr{Wagenmakers} (2014)، فصل 11.} نظریه‌ی تشخیص
  علامت به‌عنوان یک مدل شناختی (بیزی)؛ در هفته‌های 9 و 13 به آن برمی‌گردیم.
  بخش مقدمه را الان بخوانید و از کد بگذرید.
\item
  \textbf{\lr{Verguts} (2022)} و \textbf{\lr{Anderson} (2014)}: فصل‌های
  مربوط به ساختن مدل و پارامترها {[}شماره‌ی فصل‌ها بررسی شود{]}.
\item
  \emph{پیش‌زمینه درباره‌ی منحنی‌های فراموشی:} \lr{Lee} \lr{\&} \lr{Wagenmakers}
  (2014)، فصل 10 (نگهداری حافظه) از یک مدل نگهداری نمایی نزدیک استفاده
  می‌کند که آن را در هفته‌های 9\lr{--}11 به روش بیزی برازش می‌کنیم.
\end{itemize}

\emph{پیش‌نمایش هفته‌ی 3:} خطی‌سازی را کنار می‌گذاریم و مدل‌ها را مستقیماً با
کمینه کردن مربع خطا برازش می‌کنیم؛ یعنی یاد می‌گیریم الگوریتم‌های بهینه‌سازی
چطور فضای پارامتر را جستجو می‌کنند.
